% Source: source/普通高考/2011/2011天津理.pdf, page 1, question 3.
% The source PDF is vector line art; the original PNG is retained for reference.
% Nodes: 开始; a=1,i=0; i=i+1; a=i×a+1; decision a>50?; 输出 i; 结束.
% Directed edges: 开始→a=1,i=0→i=i+1→a=i×a+1→a>50?;
% 是→输出 i→结束; 否→left loop→shared stem above i=i+1→i=i+1.
\begin{tikzpicture}[
  >=Stealth,
  line width=0.55pt,
  line cap=round,
  line join=round,
  every node/.style={font=\normalsize,anchor=center,align=center,
    execute at begin node={\setlength{\parindent}{0pt}\noindent}},
  flowbox/.style={draw,minimum height=.68cm,inner xsep=4pt,inner ysep=2pt,
    anchor=center,align=center,
    execute at begin node={\setlength{\parindent}{0pt}\noindent}},
  terminal/.style={flowbox,rounded corners=.18cm,minimum width=1.35cm,
    minimum height=.72cm},
  io/.style={flowbox,shape=trapezium,trapezium left angle=72,
    trapezium right angle=108,minimum width=2.40cm,minimum height=.80cm},
  decision/.style={flowbox,shape=diamond,aspect=2.7,minimum width=4.30cm,
    minimum height=1.02cm,inner sep=1pt},
  flowarrow/.style={->,line width=0.55pt},
]
  \node[terminal] (start) at (0,10.55) {开始};
  \node[flowbox,minimum width=3.95cm] (init) at (0,9.15) {$a=1,i=0$};
  \node[flowbox,minimum width=2.75cm] (inc) at (0,7.65) {$i=i+1$};
  \node[flowbox,minimum width=4.35cm] (update) at (0,6.05)
    {$a=i\times a+1$};
  \node[decision] (test) at (0,4.25) {$a>50?$};
  \node[io] (output) at (0,2.65) {输出 $i$};
  \node[terminal] (finish) at (0,1.05) {结束};
  % Keep the shared stem clear of the init box.
  \coordinate (loopjoin) at (0,8.35);

  \draw[flowarrow] (start.south) -- (init.north);
  \draw (init.south) -- (loopjoin);
  \draw[flowarrow] (inc.south) -- (update.north);
  \draw[flowarrow] (update.south) -- (test.north);
  \draw[flowarrow] (test.south) -- node[right=2pt] {是} (output.north);
  \draw[flowarrow] (output.south) -- (finish.north);

  % 否 branch returns left and enters the shared stem above i=i+1.
  \coordinate (loop-left) at (-3.25,4.25);
  \coordinate (loop-entry) at (-3.25,8.35);
  \coordinate (loop-tip) at (-0.60,8.35);
  \draw (test.west) -- node[above=2pt] {否} (loop-left);
  \draw (loop-left) -- (loop-entry);
  \draw[flowarrow] (loop-entry) -- (loop-tip) -- (loopjoin);
  \draw[flowarrow] (loopjoin) -- (inc.north);

  \path[use as bounding box] (-3.55,0.65) rectangle (2.35,10.95);
\end{tikzpicture}
